% Section 5. Sources: 641_tape_window_bounds/PROOF.md Lemma L; 607_bounds_paper/paper.tex section 5 (finite region);
% 805_entropy_certificates/PROOF.md section 1 (Theorem F, proof steps 1-4, remarks); 808_referee_flow_bound/FINAL.md;
% 1252_referee_lower_theory/PROOF.md sections 0-2 (strip, memory m, Proposition A, Theorem F', remarks) and FINAL.md
% section 3 (tests on strips); 607_bounds_paper/paper_word.tex section 4 (exposition).
%%%%%%%%%%%%%%%%%%%%%%%%%%%%%%%%%%%%%%%%%%%%%%%%%%%%%%%%%%%%%%%%%%%%%%%%
\section{From flows to a lower bound}\label{sec:flows}
%%%%%%%%%%%%%%%%%%%%%%%%%%%%%%%%%%%%%%%%%%%%%%%%%%%%%%%%%%%%%%%%%%%%%%%%

For a lower bound we must exhibit many closed walks. An upper bound may count too much; a lower bound may count too
little. This section first recalls the weight inequality with the direction reversed, which gives lower bounds, and
then describes the entropy method used for the lower bound of Theorem~\ref{thm:main}.

\subsection{Counting inside a finite region}

\begin{lemma}\label{lem:region}
Let $\psi\ge0$ be a weight on the tapes of $\cls$ that is zero on every tape with more than $N$ letters and has
$\psi(\text{root})>0$. Suppose that $\sum_{T'}\psi(T')\ge R\,\psi(T)$ at every tape $T$ with $\psi(T)>0$, the sum running
over the successors. Then $\mu\ge R^2$.
\end{lemma}

\begin{proof}
Follow all histories from the root and delete a history when it reaches a tape of weight zero. Each move multiplies
the total weight of the survivors by at least $R$, so after $t$ moves at least $R^t\psi(\text{root})/\max\psi$
histories survive. Each survivor ends at a tape with at most $N$ letters and can be closed by Lemma~\ref{lem:sweep} in a
bounded number of moves, and different survivors with the same end tape give different closed walks. The closed walks
have boundedly many different lengths, so one length gets a fixed fraction of them. Neither the extra moves nor the
fraction changes the growth rate.
\end{proof}

The earlier lower bound $12.0535$~\cite{BoundsPaper} came from this lemma on the bracket encoding and needed $1.26$
billion table entries: the region must be large and the weight must fade toward its boundary. The entropy method
works on infinite tapes and has no such boundary.

\subsection{Counting by entropy}

Choose a random rule for the next move that looks only at the letters near the head. Intuitively, if the rule makes
on average $e^h$ equally likely choices per move, then it produces about $e^{ht}$ different histories of $t$ moves. If
the rule is stationary, so that the tape seen from the head looks statistically the same at all times, and if on
average it neither lengthens the tape nor moves the head to one side, then a typical history stays within a short
stretch of tape, and Lemma~\ref{lem:sweep} closes it at small cost. Then $c_t$ is at least about $e^{ht}$, and
\eqref{eq:fekete} gives $\mu\ge e^{2h}$. The work is to make such a rule exactly stationary on infinite tapes while
writing it with finitely many numbers. Theorem~F below makes this precise.

\subsection{Stationary flows}

\emph{The strip.} Fix $K\ge2$, or $K=\infty$. Read a tape as a path of heights: a starting height, then one step
up or down per letter. For a tape reached from the root the heights are the gap numbers \eqref{eq:path} plus
$\lfloor K/2\rfloor$; moves never change the starting height. A tape lies in the \emph{strip of height $K$} if all
heights lie in $\{0,\ldots,K-1\}$. On the strip, a split is allowed only if the new path stays in the strip; head
moves and merges never leave it. Closed walks from the root on a strip count the meanders whose labels stay in the
fixed band $-\lfloor K/2\rfloor\le g\le K-\lfloor K/2\rfloor$ of $K+1$ values (Section~\ref{sec:tape}); for even
$K$ this is $-K/2\le g\le K/2$. Their growth rate $\Lambda_K$ is at most $\mu$.

\emph{The reference law.} An infinite tape is a two-sided infinite path of heights in the strip with a head at one
letter. The \emph{reference law} $\beta$ is a stationary Markov chain on such paths with \emph{memory} $m$: the
probability of the next step depends on the current height and the last $m$ steps. We assume that every path of $m+2$
heights in the strip has positive probability and that the chain is irreducible. For a finite word $w$ with given
heights, $\beta(w)$ is its probability. The reference law is a fixed background; it describes the tape far from the
head.

\emph{Windows and flows.} Fix a radius $r$. A \emph{window} $w$ is the word of the $2r+1$ letters centered on the head,
together with its heights. A \emph{flow} is a family of numbers $J(w,\xi)\ge0$, the \emph{flow values}, one for each
window $w$ and each move $\xi$ that applies at $w$ in the strip, with total $1$. Intuitively, $J(w,\xi)$ is the
probability that the window is $w$ and the next move is $\xi$. Put
\[ \nu(w)=\sum_\xi J(w,\xi),\qquad p(\xi\mid w)=J(w,\xi)/\nu(w). \]
The flow defines a random infinite tape and a random move: the window has law $\nu$; given the window, the letters
beyond it continue as the chain $\beta$, independently on the two sides, from the $m$ letters at each end of the
window; then the move $\xi$ is chosen with probability $p(\xi\mid w)$. Intuitively, the window is chosen by the flow
and the far tape by the reference law.

\emph{The conditions.}
\begin{itemize}
\item[(S)] \emph{Stationarity.} For every word $Y$ of $2r+3$ letters, with heights, let $Y^\circ$ be its middle
$2r+1$ letters. The probability that after the move the $2r+3$ letters around the head read $Y$ equals the probability
that they read $Y$ now:
\[ \sum_{w,\xi}J(w,\xi)\,\beta_{w,\xi}(Y) \;=\; \nu(Y^\circ)\,\frac{\beta(Y)}{\beta(Y^\circ)}, \]
where $\beta_{w,\xi}(Y)$ is the conditional probability, given the window $w$ and the reference law outside it, that
the move $\xi$ produces a tape whose $2r+3$ central letters are $Y$. We call these the \emph{balance equations}.
\item[(D)] \emph{Zero drift.} $\sum J(w,\xi)\,\Delta a(\xi)=\sum J(w,\xi)\,\Delta b(\xi)=0$, with the changes
$\Delta a,\Delta b$ of the side lengths from \eqref{eq:drift}. Intuitively, on average the tape does not grow and the
head does not travel.
\item[(P)] \emph{Positivity.} $J(w,\Lm)>0$ for every window $w$. It is used in Step~2 below to prove
ergodicity.
\end{itemize}
For a fixed reference law these are finitely many linear conditions on $J$. The following theorem, which we call
Theorem~F (for flow), turns them into a lower bound.

\begin{theoremF}
Let $r\ge m$, and let $J$ be a flow satisfying (S), (D) and (P). Then
\[ \mu\;\ge\;e^{2h},\qquad h\;=\;\sum_{w,\xi}J(w,\xi)\,\log\frac{\nu(w)}{J(w,\xi)} , \]
where terms with $J(w,\xi)=0$ are left out. In particular, if $h\ge\log R$ then $\mu\ge R^2$.
\end{theoremF}

The number $h$ is the entropy of the next move given the window, averaged over the stationary process: the
uncertainty of one move.

\subsection{Proof of Theorem F, in outline}

\emph{Step 1: the balance equations give stationarity of the whole tape.} Let the infinite tape have the law of the flow,
and make one move. A move changes only the head letter and its neighbors. Given the $2r+3$ letters around the new
head, the letters beyond them are again continuations of the chain from the $m$ letters at each end: for a head move
they are the old letters, for a split the old ends survive because $r\ge m$, and for a merge the two letters newly
brought into view are fresh chain letters. So the law of the new infinite tape is determined by the law of its $2r+3$
central letters, and (S) says that this law is the old one. The process is stationary. The radius condition $r\ge m$
is needed: in tests on small strips that compare the ranks of the linear systems, the balance equations with
$r=m-1$ do not imply those at larger radius, while with $r\ge m$ they do.

\emph{Step 2: the stationary process is ergodic.} By (P) the law of the flow has a positive density with respect to
$\beta$, so the two laws have the same null sets. Let $f$ be a bounded harmonic function, that is,
$\mathbb{E}[f(\omega_1)\mid\omega_0]=f(\omega_0)$ for the process $\omega_0,\omega_1,\ldots$ of infinite tapes. By
stationarity and harmonicity,
$\mathbb{E}\bigl(f(\omega_1)-f(\omega_0)\bigr)^2=2\mathbb{E}f(\omega_0)^2-2\mathbb{E}f(\omega_0)f(\omega_1)=0$, so
$f$ is unchanged by every move of positive probability, in particular by the move $\Lm$, which is a shift of the
infinite tape. The shift is ergodic for the irreducible chain $\beta$ (its period two does not matter), so $f$ is
constant. A stationary Markov process whose bounded harmonic functions are constant is ergodic.

\emph{Step 3: counting.} Write $t$ for the number of moves. Let $\Delta^a_i$ and $\Delta^b_i$ be the total changes
of the side lengths after $i$ moves, and $I_t$ the sum of $-\log p(\xi_i\mid w_{i-1})$ over the first $t$ moves. By the
ergodic theorem and (D), $\Delta^a_i/t$ and $\Delta^b_i/t$ tend to zero uniformly in $i\le t$, and $I_t/t$ tends to
$h$. Fix $\varepsilon>0$. For large $t$, with probability at least $1/2$, $|\Delta^a_i|+|\Delta^b_i|\le\varepsilon t$
for all $i\le t$ and $I_t\ge t(h-\varepsilon)$. From some fixed infinite tape this happens with conditional
probability at least $1/2$, and each move sequence in the event has probability at most $e^{-t(h-\varepsilon)}$, so
there are at least $\frac12e^{t(h-\varepsilon)}$ of them. Cut out the letters within distance
$\lceil\varepsilon t\rceil+r+2$ of the head. All these move sequences run on this finite piece, because the head never
comes within two letters of its ends. Pad the piece to a tape of $\cls$, build it from the root and return to the root
afterwards by Lemma~\ref{lem:sweep}; both cost $O(\varepsilon t)+O(1)$ moves, in the \emph{unrestricted} machine
(without the strip). Different move sequences give different closed walks, with lengths between $t$ and
$t+O(\varepsilon t)+O(1)$, so for one length $t'$ in this range $c_{t'}\ge e^{t(h-\varepsilon)}/\mathrm{poly}(t)$. With
$c_{t'}\le\mu^{t'/2}$ from \eqref{eq:fekete}, letting $t\to\infty$ and then $\varepsilon\to0$ gives
$h\le\frac12\log\mu$.\qed

\subsection{Remarks}

Only the middle part of the closed walks is confined to the strip; the counted walks are walks of the unrestricted
machine. So the bound concerns all meanders, but it can never exceed $\Lambda_K$. For a fixed reference law the
conditions are linear in $J$ and $h$ is concave in $J$. A flow is \emph{exactly feasible} if it satisfies (S), (D)
and (P) exactly, not up to rounding; every exactly feasible flow, optimal or not, gives a bound.

The method was tested where the answer is known. On the strip of three levels the growth rate
$\Lambda_3=10.4245356612\ldots$ is known exactly. The method reproduces it at radius $1$, $2$ and $3$, and to
$15$ digits with a memory-three reference law, and in no test does it exceed it. Without the drift condition
(D) the same computation gives $10.883$, above the true value. So (D) is essential, and the test is able to detect a
wrong method. Theorem~F was proved in the project for a reference law of memory two that looks at the letters only,
with no strip, and later written out for strips and for laws of any memory $m$ with $r\ge m$; both forms were read
independently (Appendix~\ref{app:where}).
